\section{How far does the one-sided zero test transfer?}
The local applications above require no hypothesis on zeros. The global strategy faces a different question. Which prime sums can be bounded by the source's determinant construction? A test case from the incoming Frobenius report \cite{oldcubic} is the splitting field $L$ of $X^3-X-1$. The cubic is irreducible modulo $2$ and its discriminant is $-23$, so its group is $S_3$ and its quadratic subfield is $K=\QQ(\sqrt{-23})$. For an unramified prime $p$, write $I(p)=1$ if the cubic remains irreducible modulo $p$, and $0$ otherwise.

\paragraph{The missing inequality.}
With Dirichlet GRH in hand, the earlier criterion for RH for $\zeta_L$ becomes the following one-sided estimate, for each $\epsilon>0$:
\begin{equation}
        \sum_{p\le x,\ p\ne23} I(p)\log p-\frac{x}{3}
           \ \le\ C_\epsilon x^{1/2+\epsilon}.                    \label{eq:cubic}
\end{equation}
In $S_3$ the three-cycle indicator is $\mathbf1_{A_3}-\mathbf1_{\{1\}}$. Add higher unramified prime powers with this weight on Frobenius powers. The resulting $\Psi_I(x)$ differs from the prime sum by $O(\sqrt{x}\log^2(2x))$. For $\Re s>1$ the Mellin transform of $\Psi_I(x)-x/3$ against $x^{-s-1}dx$ on $[1,\infty)$ is
\[
 \frac1s\left(-\frac12\frac{\zeta'_K}{\zeta_K}
                 +\frac16\frac{\zeta'_L}{\zeta_L}+U(s)\right)
                 -\frac{1}{3(s-1)},
\]
where the omitted ramified factors contribute $U$ holomorphic for $\Re s>0$. This follows by comparing the Euler factors over each unramified prime. The pole at $1$ cancels. Dirichlet GRH gives zero-freeness of $\zeta_K$ for $\Re s>1/2$ away from its pole, so each zero $\rho$ of $\zeta_L$ to the right of the line produces a residue $\ord_\rho(\zeta_L)/(6\rho)$.

If \eqref{eq:cubic} holds, choose $\theta=1/2+\epsilon<1$ and a constant $C$ making $Cx^\theta-\Psi_I(x)+x/3$ nonnegative. Its Mellin transform has a meromorphic continuation to $\Re s>\theta$ whose possible real poles there have strictly negative residue. Landau's theorem forces a real singularity at a finite abscissa of convergence, whereas such a negative residue would make the transform negative immediately to the right. Thus the abscissa is at most $\theta$ and the integral is holomorphic for $\Re s>\theta$. Taking $\epsilon$ arbitrarily small and applying the functional equation proves the desired RH. Conversely RH for $L$ and $K$ gives \eqref{eq:cubic} by their prime ideal theorems \cite{gm}. This recovers the specific prior reduction and locates its unproved input.

\paragraph{The exact error of a fixed congruence approximation.}
Suppose one tries to express the weight $I(p)$ by a real function $w$ of $p\bmod Q$, for a fixed modulus $Q\ge1$. The best possible mean square error is
\begin{equation}
 \min_w\ \lim_{x\to\infty}\frac{1}{\#\{p\le x\}}
 \sum_{\substack{p\le x\\p\nmid23Q}}\!\big(I(p)-w(p\bmod Q)\big)^2
   =
 \begin{cases} 2/9,&23\nmid Q,\\ 1/9,&23\mid Q. \end{cases}       \label{eq:approx}
\end{equation}
Indeed, $L\cap\QQ(\zeta_Q)$ is abelian over $\QQ$, hence is either $\QQ$ or $K$. The quadratic field has conductor $23$, so the second case occurs exactly for $23\mid Q$. Chebotarev's theorem for the compositum gives uniform Frobenius distributions within the compatible fibers for conjugacy-invariant tests \cite{cheb}. If the intersection is trivial, the conditional probability of a three-cycle in every reduced class is $1/3$. Otherwise half the classes force odd parity, with probability zero of a three-cycle, and half force even parity, with probability $2/3$. In each class the minimizer is its conditional probability and the error is the Bernoulli variance. This gives \eqref{eq:approx}. Even the best binary congruence rule, with the modulus freely chosen, misclassifies a proportion $1/6$ of primes. It predicts irreducibility on even classes and reducibility on odd classes.

This calculation excludes an exact replacement of the prime weight by any fixed collection of Dirichlet characters. The error in \eqref{eq:cubic} is a signed sum and could still exhibit cancellation. Establishing that cancellation by a determinant estimate for the three-cycle weight would supply the missing analytic input.

\section{Verification and further questions}
In the four-step classification a finite enumeration is part of the proof. Two implementations reproduce the seven-prime list. One searches for a normalized witness and one counts the entire curve by convolving two-coordinate moment tables. The latter checks all 62 primes from $11$ through $317$. Separately, reduced Laplacians at $p=11,13$ give four-step valuations $46,62$, and an integer cofactor gives the six-step valuation $118$ at $p=13$.

For branch precision, a monic-remainder computation checks 1,400 combinations of contact profiles and precisions. A separate Laurent computation compares 640 cases with literal branches $(1+X)^b-1$. A third check constructs the observation matrices over $\ZZ/p^r$ and compares image sizes before and after adjoining the desired Taylor coefficients. It accepts eight sufficient cutoff vectors and detects loss of information at every one of their 25 one-coordinate reductions. A further 6,128 finite residue-ball designs test the minimax formula. These finite checks illustrate the general interpolation proofs. The new theorems are proved here on paper and remain to be formalized in Lean. Publication priority remains unverified.

\paragraph{Further research questions.}
For the branch problem, determine the least eventual period of $t_b(r)-Dr$ directly from the contact multiplicities. For graph towers, which root configurations of the quartic, sextic and higher directional moments can occur on a single plane for $m=5$? For the cubic splitting field, can a geometric construction produce \eqref{eq:cubic}? Formula \eqref{eq:approx} quantifies the information that would have to come from outside any fixed congruence-class description.

\begin{thebibliography}{99}\small
\bibitem{repo} L. Sernau and collaborators, \emph{GRH for Dirichlet $L$-functions, in Lean 4}. Shorter rewriting of the Alp\"oge--Claude proof. \href{https://github.com/lsernau/grh-lean/tree/72fead102f663d5c7a407964a13dc0cf9d977207}{Commit 72fead102f66}, 25 November 2026.
\bibitem{sutherland} A. V. Sutherland, \emph{Verification files for the shorter Lean proof of RH and of GRH for Dirichlet $L$-functions}, 26 November 2026. \href{https://math.mit.edu/~drew/shortgrh/}{Test logs and challenge statements}.
\bibitem{oldbranches} ProveIt incoming, \emph{Optimal signings of square cycles and sharp Lubin--Tate jet recovery}. Integral-precision proposition and further question 3. \href{https://github.com/VladimirReshetnikov/ProveIt/blob/1677488425e45c3df14149c804014133798136a5/docs/incoming/grh_transfers_cycle_signings_and_lubin_tate.zip}{Pinned manuscript archive}.
\bibitem{oldgraphs} ProveIt incoming, graph towers and finite rings report, \texttt{moments.tex} and \texttt{endmatter.tex}, question 2. \href{https://github.com/VladimirReshetnikov/ProveIt/blob/1677488425e45c3df14149c804014133798136a5/docs/incoming/grh_method_transfers_graph_towers_and_finite_rings.zip}{Pinned manuscript archive}.
\bibitem{effective} ProveIt incoming, \emph{When does the spanning-tree formula start? Effective Iwasawa polynomials from cyclotomic valuations}. \href{https://github.com/VladimirReshetnikov/ProveIt/blob/1677488425e45c3df14149c804014133798136a5/docs/incoming/Effective_Graph_Towers.zip}{Pinned manuscript archive}.
\bibitem{oldcubic} ProveIt incoming, McNew density and GRH transfers report, \texttt{frobenius.tex}. \href{https://github.com/VladimirReshetnikov/ProveIt/blob/1677488425e45c3df14149c804014133798136a5/docs/incoming/mcnew_density_and_grh_transfers.zip}{Pinned manuscript archive}.
\bibitem{bs} L. Berger and J. Sprang, \emph{Integer-valued polynomials and $p$-adic Fourier theory}, \S4. \href{https://arxiv.org/html/2502.18053v2}{arXiv:2502.18053v2}.
\bibitem{lt} J. Lubin and J. Tate, \emph{Formal complex multiplication in local fields}, Ann. of Math. \textbf{81} (1965), 380--387, Theorems 1--2. \href{https://doi.org/10.2307/1970622}{Publisher}.
\bibitem{cs} P. L. Clark and U. Schauz, \emph{Functional degrees and arithmetic applications, I}, Theorem 4.7 and its antecedents due to Weisman and Wilson. \href{https://arxiv.org/abs/2201.02763}{arXiv:2201.02763}.
\bibitem{dv} S. DuBose and D. Valli\`eres, \emph{On $\ZZ_\ell^d$-towers of graphs}, Algebr. Comb. \textbf{6} (2023), 1331--1346. \href{https://doi.org/10.5802/alco.304}{Publisher}.
\bibitem{dlrv} C. Dion, A. Lei, A. Ray and D. Valli\`eres, \emph{On the distribution of Iwasawa invariants associated to multigraphs}, Nagoya Math. J. \textbf{253} (2024), 48--90. \href{https://doi.org/10.1017/nmj.2023.18}{Publisher}.
\bibitem{weil} A. Weil, \emph{Sur les courbes alg\'ebriques et les vari\'et\'es qui s'en d\'eduisent}, Hermann, 1948. \href{https://denisevellachemla.eu/Weil-courbes-varietes.pdf}{Monograph}.
\bibitem{gm} L. Greni\'e and G. Molteni, \emph{Explicit versions of the prime ideal theorem for Dedekind zeta functions under GRH}, Math. Comp. \textbf{85} (2016), 889--906. \href{https://air.unimi.it/handle/2434/448115}{Article record}.
\bibitem{cheb} J. C. Lagarias, H. L. Montgomery and A. M. Odlyzko, \emph{A bound for the least prime ideal in the Chebotarev density theorem}, Invent. Math. \textbf{54} (1979), 271--296. \href{https://doi.org/10.1007/BF01390234}{Publisher}.
\end{thebibliography}
